\documentclass[10pt]{article}

\usepackage[margin=0.7in]{geometry}
\usepackage{booktabs}
\usepackage{array}
\usepackage{hyperref}
\usepackage{enumitem}
\setlist{nosep, leftmargin=1.4em}
\setlength{\parskip}{0.3em}
\setlength{\parindent}{0pt}
\pagestyle{empty}

\title{\vspace{-2em}\textbf{\large VAKRA-Advanced: Tool-Using LLM Agents for Multi-Hop,
Multi-Source Reasoning}\\[0.2em] \normalsize Phase 1 Report --- Environment, Baseline, and
Diagnostics}
\author{\small Anurag Kaushal (\texttt{[Roll No.]}) \quad Aayush Pandey (\texttt{[Roll No.]})
\quad Piyush Agarwal (\texttt{[Roll No.]})}
\date{\small October 30, 2026}

\begin{document}
\maketitle
\vspace{-1.5em}

\textbf{Task.} VAKRA-Advanced (SemEval 2027) is about building an agent that answers a user's
question by using external tools, rather than answering from what a language model already
"knows". Each domain (e.g.\ hockey, sales, movies) comes with its own set of tools -- anywhere
from 6 to over 300 -- and the agent must pick the right one, fill in its inputs correctly, and,
when one tool's output is needed as input to another, chain several tool calls together.
Sometimes the tools alone are not enough, and the agent must also look up supporting text from a
document collection. Some questions additionally come with a plain-English rule (a "policy")
saying which tools may or may not be used, and the agent must follow it, even if that means
refusing to answer. The benchmark covers four such capabilities: Tool Selection, Multi-hop API
Reasoning, API Chaining, and Multi-source reasoning with policies. To check an answer, the
official evaluator re-runs the agent's exact tool calls against the real tools and checks that
they match the expected ones; if not, an LLM judge decides whether the final answer is still
correct and properly backed up by the evidence the agent actually collected. Only the last
question of a conversation is graded, and a correct, honest "I cannot answer" (when the policy
genuinely blocks every useful tool) counts as fully correct. Following the proposal, our project
takes up Tool Selection and Multi-hop API Reasoning first, the multi-source/policy capability if
compute allows, and leaves API Chaining for later.

\textbf{Phase 1 work.} Phase 1 was about setting everything up correctly and measuring where
things stand today, before changing anything. We got the official benchmark running-- its databases and
tools included. We
then fixed a set of evaluation groups, held out by domain so that results are never just
memorizing seen data, including one representative slice, \texttt{dev-eval}, used for every
number below so that all comparisons are on exactly the same questions. On this slice, we ran the
benchmark's own ready-made agent, unchanged, as our baseline, once with the full list of
candidate tools per question (up to 128, the benchmark's default). 

\textbf{Scale and cost.} The benchmark spans 1,300--1,900 training questions per capability
across 28--40 domains, with tool catalogs of 6--328 tools per domain (median $\sim$105). Tool
Selection is always a single tool call; Multi-hop questions are almost always exactly two hops
(90\%); API Chaining questions often need 3 or more chained calls (60\%); and multi-source
conversations rarely need document retrieval at all ($\sim$3--7\% of turns), always alongside an
API call. Running one baseline configuration on \texttt{dev-eval}, agent and judge calls
included, cost about \$1--2 in language-model usage.

\begin{table}[h!]
\centering\small
\vspace{-0.3em}
\begin{tabular}{@{}lrrrr@{}}
\toprule
Configuration & Cap 2 ($n{=}120$) & Cap 3 ($n{=}35$) & Cap 4 ($n{=}60$) & Macro avg.\ (2--4) \\
\midrule
Default shortlist (128) & 0.408 & 0.343 & 0.350 & 0.367 \\
Reduced shortlist (32)  & 0.425 & 0.171 & --- & --- \\
\bottomrule
\end{tabular}
\caption{\small Baseline score on \texttt{dev-eval}. 95\% CI half-widths $\approx\pm0.09$ (Cap 2),
$\pm0.16$ (Cap 3), $\pm0.12$ (Cap 4).}
\vspace{-0.5em}
\end{table}

\textbf{Findings.}
\begin{itemize}
\item \textbf{Retrieval, not reasoning, bounds Capability 3.} With a shortlist of 32 tools, the
correct tool set is present for only 71\% of Capability 3 questions, versus 94\% with 128 tools;
a combined (lexical + embedding) retrieval method recovers part of this, reaching 80\% at the
same shortlist size. This gap alone explains why the 32-tool baseline score is roughly half the
128-tool score on Capability 3, while Capability 2 is barely affected.
\item \textbf{About a fifth of Capability 2--3 questions fail before reasoning even matters.} The
agent's tool call is missing a required input, and the service rejects the call outright instead
of giving the agent a chance to fix it.
\item \textbf{7--23\% of questions reach the right tool results or the right answer but are still
marked wrong}, because the answer is judged as not properly backed up by that evidence.
\item \textbf{Deliberate refusal is a missed opportunity on Capability 4.} About 18\% of
\texttt{dev-eval} conversations are decided purely by whether the agent correctly says "I cannot
answer"; the baseline agent never does this on purpose.
\item \textbf{About 17\% of Capability 4 conversations are simply too long} for the language
model to process at all, regardless of how many tools are offered.
\end{itemize}

\end{document}
